%% kapa: exam=1 round=29 year=2018 period=2 form=A subject=appraisal_law subject_ko=감정평가관계법규 seq=35 number=35
%% source: Q-Net 감정평가사 기출문제 (cst00309) artlSeq=5202808 file=extracted/1008/5202808/2018년 제29회 감정평가사 1차시험 기출문제/2018년 제29회 감정평가사 1차시험 기출문제/2교시 A형.pdf
%% conversion: kordoc 4.24.51 -> tools/gen1.py (auto); review=proofread-sonnet; flags=none
\begin{question}{35}
부동산등기법상 건물의 표시에 관한 등기에 대한 설명으로 옳지 않은 것은?
\begin{choices}
  \item 등기관이 대지권등기를 하였을 때에는 직권으로 대지권의 목적인 토지의 등기기록에 소유권, 지상권, 전세권 또는 임차권이 대지권이라는 뜻을 기록하여야 한다.
  \item 건물이 멸실된 경우에는 그 건물 소유권의 등기명의인은 그 사실이 있는 때부터 2개월 이내에 그 등기를 신청하여야 한다.
  \item 존재하지 아니하는 건물에 대한 등기가 있을 때에는 그 소유권의 등기명의인은 지체 없이 그 건물의 멸실등기를 신청하여야 한다.
  \item 구분건물로서 그 대지권의 변경이나 소멸이 있는 경우에는 구분건물의 소유권의 등기명의인은 1동의 건물에 속하는 다른 구분건물의 소유권의 등기명의인을 대위하여 그 등기를 신청할 수 있다.
  \item 1동의 건물에 속하는 구분건물 중 일부만에 관하여 소유권보존등기를 신청하는 경우에는 나머지 구분건물의 표시에 관한 등기를 동시에 신청하여야 한다.
\end{choices}
\end{question}
\answer{2}
